\section{A surface criterion for the switch}\label{sec:surface}

To prove Theorem~\ref{thm:switch}, we will construct four degree-one
classes on one surface with a nonzero cup product.  The criterion
below explains why this is sufficient.  The four maps give an
algebraic surface class that tests the switch line nontrivially; four
divisor classes then convert that test into a vector on the line
itself.

\begin{proposition}[Surface criterion]\label{prop:surface-switch}
Let \(E\) be a Galois CM field with conjugation \(c\), and let
\(v_1,v_2,v_3,v_4:G_E\to\{1,-1\}\) be odd sign functions satisfying
\(v_1+v_2=v_3+v_4\).  Put \(\lambda=(1,1,c,c)\).
Suppose that \(S\) is a smooth connected projective complex surface
and that there are rational Hodge maps
\[
 h_i:U(v_i)\longrightarrow H^1(S,\Q),\qquad 1\leq i\leq4,
\]
and nonzero vectors
\(\alpha_i\in U(v_i)_{\lambda_i}\subset U(v_i)_{\Qbar}\) such that
\begin{equation}\label{eq:surface-period}
 \int_S h_1(\alpha_1)\smile h_2(\alpha_2)
             \smile h_3(\alpha_3)\smile h_4(\alpha_4)\ne0.
\end{equation}
Here each \(h_i\) is extended to \(\Qbar\)-coefficients.  Then the
line in Equation~\eqref{eq:switch-line} is algebraic.
\end{proposition}

\begin{proof}
By Lemma~\ref{lem:weight-one-maps}, there are positive integers \(m_i\)
and morphisms \(f_i:S\to B(v_i)\) with \(f_i^*=m_i h_i\).  Therefore
\begin{equation}\label{eq:surface-morphism-period}
 k:=\int_S f_1^*\alpha_1\smile f_2^*\alpha_2
              \smile f_3^*\alpha_3\smile f_4^*\alpha_4
 \in\Qbar^\times.
\end{equation}
Nonvanishing follows from Equation~\eqref{eq:surface-period}, since
the new value is multiplied by \(\prod_i m_i\).  The value is
algebraic because the pullbacks, cup product, and evaluation on
\([S]\) are rational operations on Betti cohomology and the
\(\alpha_i\) have algebraic Betti coefficients.

Put \(A=\prod_iB(v_i)\), \(n=\dim A=4s\), and \(f=(f_i):S\to A\).
Let
\[
 z=f_*[S]\in\CH^{n-2}(A),\qquad
 Z=\operatorname{cl}(z)\in\mathcal Z^{2n-4}(A).
\]
For \(t_\alpha=\alpha_1\otimes\cdots\otimes\alpha_4\) in the
four-\(H^1\) Künneth summand of \(H^4(A,\Qbar)\), the projection
formula gives
\begin{equation}\label{eq:surface-functional}
 \int_A Z\smile t_\alpha=k.
\end{equation}
Thus \(Z\) gives a nonzero functional on the switch line; it is not
itself a vector on that line.  In particular \(f(S)\) cannot have
dimension less than two, since then the proper pushforward of the
two-cycle \([S]\) and the displayed pairing would be zero.

For each \(i\), take a copy \(B(v_i)'\) with its transported
\(E\)-action, and choose a nonzero vector
\(\beta_i\in U(v_i)'_{c\lambda_i}\) with algebraic Betti coefficients.
Under the Künneth inclusion define the pure class
\[
 D_i=\alpha_i\otimes\beta_i
 \in H^2\bigl(B(v_i)\times B(v_i)',\Qbar\bigr).
\]
If \(\sigma\in\Gal(\Qbar/\Q)\) restricts to \(g\in G_E\), then
Equation~\eqref{eq:cm-label-conjugation} puts its two factors on the
lines labeled
\[
 g\lambda_i,\qquad gc\lambda_i=cg\lambda_i.
\]
These have opposite Hodge types.  This uses centrality of \(c\) in
\(G_E\) only, and remains valid when the scalar coordinates of
\(\alpha_i\) and \(\beta_i\) lie outside \(E\).  Every scalar
conjugate of \(D_i\) therefore has type \((1,1)\).  Scalar descent
and the Lefschetz \((1,1)\) theorem give
\begin{equation}\label{eq:surface-divisor-kernels}
 D_i\in
 \mathcal Z^2\bigl(B(v_i)\times B(v_i)'\bigr)\otimes_{\Q}\Qbar.
\end{equation}
The entire cohomology class \(D_i\) is algebraic here.  It was defined
in the pure \(H^1\otimes H^1\) summand before applying descent, so
there is no need to project an arbitrary divisor class onto that
summand.

Set \(A'=\prod_iB(v_i)'\).  Let \(p:A\times A'\to A\) and
\(p':A\times A'\to A'\) be the projections, and let
\(q_i:A\times A'\to B(v_i)\times B(v_i)'\) select the \(i\)-th old
and new slots.  Write
\[
 K=\prod_{i=1}^4q_i^*D_i,\qquad
 t_\beta=\beta_1\otimes\cdots\otimes\beta_4.
\]
Order the old factors first and the new factors second in the
Künneth identification for \(A\times A'\).  Moving the new
degree-one factor of \(D_i\) past the later old degree-one factors
makes \(3+2+1=6\) sign changes.  Hence
\[
 K=(-1)^6\,t_\alpha\otimes t_\beta=t_\alpha\otimes t_\beta.
\]
In particular \(K\) has degree four in the old factors and degree
four in the new factors.  Integration over \(A\) consequently gives
the exact formula
\begin{equation}\label{eq:surface-to-tensor}
 p'_*\bigl(p^*Z\smile K\bigr)
   =\left(\int_A Z\smile t_\alpha\right)t_\beta
   =k\,t_\beta.
\end{equation}
Any other Künneth component of \(Z\) contributes zero to this
integration.  The codimension of the left side is
\((n-2)+4-n=2\), equivalently its cohomological degree is
\((2n-4)+8-2n=4\).

By Equation~\eqref{eq:surface-divisor-kernels}, the left side of
Equation~\eqref{eq:surface-to-tensor} is algebraic with
\(\Qbar\)-coefficients.  Since \(k\in\Qbar^\times\), division by \(k\)
proves algebraicity of the line generated by \(t_\beta\), whose
labels are \((c,c,1,1)\).  Ordinary scalar complex conjugation fixes
the rational cycle span on the same \(A'\) and changes these labels
to \((1,1,c,c)\).  Identifying the copied factors with the original
ones gives the asserted line on \(A\).
\end{proof}

The theta construction produces complex classes, and the arithmetic
argument will place them in spans of rational Hodge images.  The
following consequence applies to those classes.

\begin{corollary}[Surface criterion for complex spans]
\label{cor:surface-span}
Let \(E\), \(v_1,\ldots,v_4\), and \(\lambda=(1,1,c,c)\) be as in
Proposition~\ref{prop:surface-switch}, and let \(S\) be a smooth
connected projective complex surface.  Suppose that, for each
\(1\leq i\leq4\), a class \(\eta_i\in H^1(S,\C)\) belongs to
\[
 \operatorname{span}_{\C}
 \{h(U(v_i)_{\lambda_i}):
       h\in\Hom_{\mathrm{Hdg}}(U(v_i),H^1(S,\Q))\}
\]
and that
\[
 \int_S\eta_1\smile\eta_2\smile\eta_3\smile\eta_4\ne0.
\]
Then the line in Equation~\eqref{eq:switch-line} is algebraic.
\end{corollary}

\begin{proof}
Fix a nonzero algebraic Betti generator of each selected line.
Each \(\eta_i\) is a finite complex linear combination of its images
under rational Hodge maps.  Expanding the nonzero cup product selects
a nonzero term of the form in Equation~\eqref{eq:surface-period}.
Thus Proposition~\ref{prop:surface-switch} applies.  This argument
uses only finite sums; all four maps have the same target
\(H^1(S,\Q)\).
\end{proof}
